\section{模型在分化：通用、Coding+Agentic、多模态交互}

导读提出“分化”和“产品”两个关键词，本章先展开第一条，核心问题是：头部模型公司不再只沿同一条 benchmark 曲线竞争，而是在能力栈和产品路径上开始分叉。硅谷模型公司这个季度开始显著分化。Google Gemini 和 OpenAI 仍然做通用能力；Anthropic 分化到 Coding 和 Agentic 模型能力；Thinking Machines Lab 分化到多模态和下一代交互；Grok 还在摸索生态位置；Meta 的原创 0-1 基因被认为较弱。最领先的几家公司像 F1 竞赛，速度极快，路线也开始不同。

\begin{figure}[H]
\centering
\includegraphics[width=0.96\textwidth]{model-company-divergence.png}
\caption{模型公司开始分化：Gemini/OpenAI 通用，Anthropic Coding+Agentic，Thinking Machines 多模态交互。自制概念图，依据 00:03:54--00:21:37 对谈内容整理。}
\end{figure}

\begin{knowledgebox}{读图：分化不是掉队，而是选择赛道}
通用模型要保持全能力领先；Coding+Agentic 下注高价值工作流；多模态/交互下注下一代入口；生态型公司需要找到与自有平台的结合点。分化说明行业从单纯追 benchmark 进入战略选择。
\end{knowledgebox}

\begin{knowledgebox}{术语消化：公司路线地图}
\begin{tabularx}{\textwidth}{p{0.22\textwidth}p{0.34\textwidth}X}
\toprule
公司/路线 & 本期定位 & 关键风险 \\
\midrule
Gemini/OpenAI & 通用模型和全能力入口 & 成本高，产品线复杂。 \\
Anthropic & Coding + Agentic 能力 & 需把技术壁垒转成更大产品入口。 \\
Thinking Machines & 多模态原生和下一代交互 & 方向新但商业路径未定。 \\
Grok/xAI & 生态位置仍在摸索 & 平台协同和品牌心智仍需证明。 \\
Meta & 开源和社交生态强，但原创 0-1 被质疑 & 模型路线和产品入口之间不稳定。 \\
\bottomrule
\end{tabularx}
\end{knowledgebox}

\begin{knowledgebox}{分化背后的能力栈}
\begin{tabularx}{\textwidth}{p{0.22\textwidth}p{0.34\textwidth}X}
\toprule
能力栈 & 代表公司/方向 & 关键资产 \\
\midrule
通用智能 & Gemini、OpenAI & 多模态、推理、搜索、产品入口和算力。 \\
Coding/Agentic & Anthropic、Claude Code & 代码环境、长程任务、企业开发者心智。 \\
多模态交互 & Thinking Machines、Gemini & 原生多模态、下一代 interface 和硬件入口。 \\
生态型模型 & Grok、Meta & 社交平台、开源生态、分发渠道。 \\
\bottomrule
\end{tabularx}
\end{knowledgebox}

\begin{knowledgebox}{模型公司战略对照表}
\begin{tabularx}{\textwidth}{p{0.18\textwidth}p{0.30\textwidth}p{0.40\textwidth}}
\toprule
公司/路线 & 当前强项 & 需要继续证明 \\
\midrule
OpenAI & ChatGPT 入口、通用能力、产品化 & 广告/商业化是否不伤用户体验。 \\
Google & TPU、搜索、广告、Workspace、Gemini & 能否把资产重新组合成默认 AI 入口。 \\
Anthropic & Claude Code、企业开发者、Agentic 能力 & 能否从 coding 扩展到更大工作流。 \\
Thinking Machines & 多模态原生、下一代交互想象 & 能否把研究路线变成强产品。 \\
Meta/xAI & 平台、开源、社交或实时数据资产 & 能否做出原创 0-1 产品心智。 \\
\bottomrule
\end{tabularx}
\end{knowledgebox}

\subsection{分化给创业者的启发}

前面给出公司地图，本节把它翻译成创业者视角。如果头部公司路线开始分化，创业者不能再只问“谁模型最强”，而要问“我站在哪条能力链上”。做 coding 工具，要看 Anthropic/Claude Code；做搜索和研究，要看 Deep Research；做多模态交互，要看 Thinking Machines 和 Gemini；做消费入口，要看 ChatGPT 的全家桶。

\begin{warningbox}{不要把分化误读为小公司机会自动变多}
分化意味着头部公司开始主动选择战场，不是放弃战场。创业者的机会在于更快、更窄、更深，而不是在巨头已经明确下场的正面战场硬拼。
\end{warningbox}

\subsection{本章小结}

模型公司分化说明行业进入策略期。通用、Coding+Agentic、多模态和生态路线都可能产生赢家，但每条路都要求不同产品、数据和商业能力。

